% ================================================================
% frame2-tex.tex
% The Tree and the Shadow — Al-Hindara Framework
% Revised Research Manuscript
% Author: Eng. Abdulla M. N. Al-Jabri — Sana'a, Yemen
% Date: October 2026
% ================================================================

\documentclass[11pt,a4paper]{article}

% ------------------------------------------------
% Packages
% ------------------------------------------------
\usepackage[utf8]{inputenc}
\usepackage[T1]{fontenc}
\usepackage{amsmath,amssymb,amsthm}
\usepackage{geometry}
\usepackage{graphicx}
\usepackage{hyperref}
\usepackage{booktabs}
\usepackage{xcolor}
\usepackage{float}
\usepackage{titlesec}
\usepackage{fancyhdr}

% ------------------------------------------------
% Page layout and metadata
% ------------------------------------------------
\geometry{margin=1in}

\hypersetup{
    colorlinks=true,
    linkcolor=blue,
    citecolor=blue,
    filecolor=magenta,
    urlcolor=cyan,
    pdftitle={The Tree and the Shadow: Reciprocal Silver Ratio and Exploratory Analysis},
    pdfauthor={Eng. Abdulla M. N. Al-Jabri}
}

% ------------------------------------------------
% Header and Footer Customization
% ------------------------------------------------
\pagestyle{fancy}
\fancyhf{}
\fancyhead[L]{\small \textbf{Eng. Abdulla M. N. Al-Jabri} -- Frame 2}
\fancyhead[R]{\small Al-Hindara Tree-Shadow Study}
\fancyfoot[C]{\thepage}
\fancyfoot[L]{\small\textit{The Tree and the Shadow}}
\fancyfoot[R]{\small \texttt{jabri62018@gmail.com}}
\renewcommand{\headrulewidth}{0.4pt}
\renewcommand{\footrulewidth}{0.4pt}
\setlength{\headheight}{14pt}

% ------------------------------------------------
% Title
% ------------------------------------------------
\title{
    \textbf{The Tree and the Shadow:}\\
    \Large Geometric Derivation of the Reciprocal Silver Ratio
    \(\sqrt{2}-1\) and an Exploratory Statistical Analysis
    of Tree Crown Ratios\\[0.5em]
    \large A Pilot Study Bridging Geometry, Biology,
    and Number Theory
}

\author{
    \textbf{Eng. Abdulla Mohammed Nasser Al-Jabri}\\
    \small Independent Researcher\\
    \small Sana'a, Yemen\\
    \small\texttt{jabri62018@gmail.com}\\
    \small ORCID:
    \href{https://orcid.org/0009-0003-3319-3822}
    {0009-0003-3319-3822}
}

\date{October 2026}

\begin{document}

\maketitle
\thispagestyle{fancy}

% ================================================================
% ABSTRACT
% ================================================================
\begin{abstract}
\noindent
This paper presents the Al-Hindara framework as a conceptual
approach to connecting geometric modeling, management, and
empirical observation. It derives the exact trigonometric identity
\[
    \tan\left(\frac{\pi}{8}\right)=\sqrt{2}-1,
\]
the reciprocal of the conventional Silver Ratio
\(1+\sqrt{2}\). The paper then explores whether this geometric
reference value is consistent with the crown-to-total-height
ratio of five mango trees reported to have been measured in
Sho'ub, Sana'a, Yemen.

The reported pilot sample has a mean crown ratio of approximately
\(0.41297\). A Two One-Sided Tests (TOST) procedure is described
using an equivalence margin of \(\pm 0.03\) around the theoretical
reference value. Under the assumptions of the test, the reported
results indicate that the sample mean falls within this
prespecified equivalence region. The small sample size and
single-location sampling limit the interpretation: these findings
do not establish a universal biological law or a causal
relationship between solar geometry and tree crown structure.

The paper also introduces a conjectural numerical modeling
proposal involving the integer \(21\), obtained by taking the floor
of the ordinate of the second positive nontrivial zero of the
Riemann zeta function. The proposed numerical connection is
exploratory and is not presented as a theorem about the zeta
function.

The paper distinguishes an exact geometric identity from
empirical observations and conjectural numerical constructions.
Its principal aim is to provide a transparent basis for further
mathematical and field investigation.

\vspace{1em}
\noindent\textbf{Keywords:}
Reciprocal Silver Ratio, Half-Angle Identity, Crown Ratio,
TOST Equivalence Testing, Mango Trees, Riemann Zeta Function,
Al-Hindara.
\end{abstract}

% ================================================================
% 1. INTRODUCTION
% ================================================================
\section{Introduction}

The Al-Hindara framework introduces the conceptual relation
\begin{equation}
    Z+C+A=1,
    \label{eq:hindara}
\end{equation}
where \(Z\), \(C\), and \(A\) are intended to represent normalized
components associated with geometry, management, and synthesis.

In this paper, Equation~\eqref{eq:hindara} is treated as a
conceptual normalization proposal. The variables must be
operationally defined and normalized before the equation can be
tested as a quantitative empirical relationship. It is not a
consequence of the Pythagorean theorem, and the present study
does not establish it as a universal law.

The paper has three objectives. First, it derives the exact
trigonometric value associated with a \(22.5^\circ\) angle measured
from the vertical. Second, it compares that value with the
reported crown-to-total-height ratios of five mango trees.
Third, it outlines a conjectural numerical construction involving
an integer derived from a nontrivial zero ordinate of the
Riemann zeta function.

These objectives have different evidential status. The
trigonometric identity is an exact mathematical result. The
tree measurements constitute a small exploratory dataset.
The proposed number-theoretic connection is a conjecture
that requires a separate mathematical derivation.

% ================================================================
% 2. GEOMETRIC DERIVATION
% ================================================================
\section{Geometric Derivation of the Reciprocal Silver Ratio}

Consider an idealized right triangle in which \(Z\) denotes the
vertical height, \(C\) the horizontal shadow length, and \(A\) the
length of the slanted ray. The Pythagorean theorem gives
\begin{equation}
    Z^2+C^2=A^2.
    \label{eq:pythagoras}
\end{equation}

Let \(\theta\) be the angle between the vertical direction and
the slanted ray. Under this convention, the side opposite
\(\theta\) is \(C\), and the adjacent side is \(Z\). Therefore,
\begin{equation}
    \tan(\theta)=\frac{C}{Z}.
    \label{eq:shadow_ratio}
\end{equation}

This angle convention differs from the conventional solar
elevation angle measured from the horizontal. If
\(\theta=22.5^\circ\) from the vertical, the corresponding
elevation angle from the horizontal is \(67.5^\circ\).

To evaluate the ratio at \(\theta=22.5^\circ\), use the
half-angle identity
\begin{equation}
    \tan\left(\frac{\phi}{2}\right)
    =\frac{1-\cos\phi}{\sin\phi}.
    \label{eq:half_angle}
\end{equation}

Setting \(\phi=45^\circ\) yields
\begin{align}
    \tan(22.5^\circ)
    &=\frac{1-\frac{\sqrt{2}}{2}}
            {\frac{\sqrt{2}}{2}}\\
    &=\frac{2-\sqrt{2}}{\sqrt{2}}\\
    &=\sqrt{2}-1\\
    &\approx 0.4142135624.
    \label{eq:reciprocal_silver}
\end{align}

The conventional Silver Ratio is \(1+\sqrt{2}\). Since
\[
    \frac{1}{1+\sqrt{2}}=\sqrt{2}-1,
\]
the value in Equation~\eqref{eq:reciprocal_silver} is its
reciprocal.

This result is exact for the specified idealized geometry.
It does not imply that tree crown structure must follow the
same ratio. Any application to biological morphology is an
empirical hypothesis that requires independent observation.

% ================================================================
% 3. EXPLORATORY FIELD STUDY
% ================================================================
\section{Exploratory Field Measurement Protocol}

\subsection{Study objective and operational definition}

The exploratory analysis compares the geometric reference value
with a crown ratio defined by
\begin{equation}
    R=\frac{H-h}{H},
    \label{eq:crown_ratio}
\end{equation}
where \(H\) is the total tree height and \(h\) is the height of
the first main branching point, both measured from the same
ground-level reference.

The variable \(R\) is dimensionless. It represents the fraction
of total height above the specified branching point under this
operational definition. It should not automatically be
interpreted as a complete measure of crown volume, crown width,
or foliage density.

The target reference is
\begin{equation}
    R_0=\sqrt{2}-1\approx0.4142135624.
    \label{eq:target}
\end{equation}

\subsection{Reported study setting and measurements}

The supplied study record describes a pilot sample of five
mango trees (\textit{Mangifera indica}) in Sho'ub, Sana'a,
Yemen, on October 14, 2026.

The reported measurement procedure specifies a metal measuring
tape with a stated precision of \(\pm1\) mm and an extendable
rigid pole used to align the measurement with the highest
leaf apex. The first main branching point is defined as the
first clear bifurcation at which each of the two branches
reaches at least \(30\%\) of the original trunk diameter.

These procedural details should be retained as a description
of the actual study only if they accurately reflect the
measurements performed. The nominal precision of the tape
does not, by itself, establish the total accuracy of a
7--9 m tree-height measurement. Alignment, access to the
highest leaf, ground slope, branch identification, and
repeatability may contribute additional uncertainty.

For reproducibility, a future study should document the
measurement instrument, the positioning procedure, the
number of repeat measurements, the treatment of sloping
ground, and the uncertainty associated with both \(H\) and \(h\).

\subsection{Raw pilot data}

Table~\ref{tab:raw_data} reproduces the five observations
provided for this analysis. The ratios are calculated from
the supplied heights using Equation~\eqref{eq:crown_ratio}.

\begin{table}[H]
\centering
\small
\begin{tabular}{ccccc}
\toprule
Tree & \(H\) (m) & \(h\) (m) & \(R=(H-h)/H\) & Reference\\
\midrule
1 & 7.80 & 4.50 & 0.42308 & \(\sqrt{2}-1\)\\
2 & 8.50 & 5.00 & 0.41176 & \(\sqrt{2}-1\)\\
3 & 7.20 & 4.20 & 0.41667 & \(\sqrt{2}-1\)\\
4 & 9.00 & 5.40 & 0.40000 & \(\sqrt{2}-1\)\\
5 & 7.50 & 4.40 & 0.41333 & \(\sqrt{2}-1\)\\
\bottomrule
\end{tabular}
\caption{Reported pilot measurements and crown ratios.
The theoretical reference is
\(R_0=\sqrt{2}-1\approx0.414214\).}
\label{tab:raw_data}
\end{table}

% ================================================================
% 4. STATISTICAL ANALYSIS
% ================================================================
\section{Statistical Analysis: TOST Equivalence Procedure}

\subsection{Hypotheses and equivalence margin}

The analysis uses the Two One-Sided Tests (TOST) procedure
for equivalence testing, following Schuirmann~\cite{schuirmann1987}.
The stated equivalence margin is
\begin{equation}
    \Delta=0.03,
\end{equation}
so the equivalence interval around the theoretical reference
is
\begin{equation}
    [R_0-\Delta,\ R_0+\Delta]
    \approx[0.384214,\ 0.444214].
    \label{eq:equivalence_bounds}
\end{equation}

Let \(\mu\) denote the population mean under the sampling model.
The TOST hypotheses are
\begin{align}
    H_{01}&:\mu\leq R_0-\Delta,\\
    H_{02}&:\mu\geq R_0+\Delta,
\end{align}
against the alternative that the mean lies strictly within
the equivalence limits.

At significance level \(\alpha=0.05\), equivalence is concluded
only when both one-sided null hypotheses are rejected.
For the conventional paired interpretation of the TOST and
the corresponding confidence-interval method, the 90\%
confidence interval for the mean must lie entirely within
the equivalence bounds.

The equivalence margin is a substantive tolerance, not a
quantity determined by the observed data. Its scientific
justification should be specified independently of the
observed results. If the margin was selected after inspecting
the sample, it should be described as exploratory rather than
prespecified.

\subsection{Descriptive statistics}

Using the five ratios in Table~\ref{tab:raw_data}, the reported
descriptive statistics are approximately
\begin{align*}
    n&=5,\\
    \bar{R}&=0.41297,\\
    s&=0.00845,\\
    SE&=0.00378.
\end{align*}

Under a one-sample Student \(t\) model with \(df=n-1=4\),
the reported 90\% confidence interval is approximately
\begin{equation}
    [0.40491,\ 0.42102].
    \label{eq:ci}
\end{equation}

This interval lies within the equivalence bounds in
Equation~\eqref{eq:equivalence_bounds}. The reported TOST
p-value is approximately \(0.000801\), using the larger of
the two one-sided p-values as the combined TOST p-value.

The numerical result should be interpreted conditional on
the sampling design and statistical assumptions. In
particular, the analysis assumes that observations are
independent and that a one-sample \(t\) procedure is
appropriate for the population mean. With only five
observations, distributional assumptions and outlier
sensitivity cannot be assessed reliably.

\subsection{Interpretation and limitations}

The TOST result indicates that the sample mean is
statistically equivalent to the theoretical reference
within the stated margin under the assumptions of the test.
It does not demonstrate that individual trees have this
ratio, that the ratio is biologically optimal, or that solar
geometry causes the observed morphology.

The small sample size and single-location sampling make
this a preliminary exploratory analysis. A stronger follow-up
study should use a clearly defined sampling frame, more trees
from multiple locations, repeatability checks, an
independently justified equivalence margin, and an
\textit{a priori} sample-size or power analysis.

If several trees share a site or management history, the
analysis should also consider whether observations are
clustered rather than independent.

% ================================================================
% 5. CONJECTURAL NUMERICAL MODEL
% ================================================================
\section{A Conjectural Numerical Construction Involving
the Riemann Zeta Function}

This section presents a numerical modeling proposal rather
than a proven connection between tree geometry and the
Riemann zeta function.

Let \(\zeta(s)\) denote the Riemann zeta function. Write
\(\gamma_2\approx21.022039\) for the ordinate of the second
positive nontrivial zero on the critical line, where zeros
are conventionally expressed as
\(s=\frac12+i\gamma\) and ordered by increasing positive
ordinate. The first positive ordinate is approximately
\(14.134725\), and the second is approximately \(21.022039\)
\cite{edwards1974}.

Define the integer parameter
\begin{equation}
    x_p=\lfloor\gamma_2\rfloor=21.
    \label{eq:xp}
\end{equation}

Within the proposed numerical model, the parameter \(x_p\)
is used in the expression
\begin{equation}
    Z(x)=x^{-5}\ln(x)\sin\left(\frac{2\pi}{x}\right)
    e^{-x/x_p},
    \qquad x>1.
    \label{eq:z_field}
\end{equation}

Equation~\eqref{eq:z_field} defines a model expression.
The present paper does not derive this expression from
properties of \(\zeta(s)\), nor does it establish that it
describes a physical field.

The parameter \(x_p=21\) is referred to within the
Al-Hindara project as the \textit{Yemeni Constant}. This
is a project-specific name for the chosen integer parameter,
not a recognized mathematical constant established by
the results of this paper.

The model also considers the numerical quantity
\begin{equation}
    K=2\pi e^{-1/42}\approx6.1353.
    \label{eq:model_coefficient}
\end{equation}

The expression in Equation~\eqref{eq:model_coefficient}
is mathematically well-defined, but its occurrence alone
does not establish a connection to the asymptotic behavior
of zeta zeros or to the Riemann critical line.

In particular, the proposed location \(2x_p=42\) for a peak
of an accumulation function cannot be verified here because
that function has not been explicitly defined and analyzed.
Establishing such a claim would require a precise definition
of the accumulation function, a domain, appropriate
regularity conditions, and a derivation locating and
classifying its extrema.

Accordingly, the recurrence of the integers \(21\) and \(42\)
is presented as a conjectural feature of the proposed
numerical construction. No theorem linking these integers
to the geometric ratio \(\sqrt{2}-1\) or to the distribution
of nontrivial zeta zeros is claimed.

% ================================================================
% 6. DISCUSSION
% ================================================================
\section{Discussion}

The results of this paper belong to three distinct categories.

First, the half-angle identity in
Equation~\eqref{eq:reciprocal_silver} is exact and follows
from standard trigonometry.

Second, the five reported crown ratios produce a sample mean
close to the geometric reference value. The equivalence
analysis quantifies this closeness relative to the chosen
margin, but it cannot establish a universal biological
pattern from such a small sample.

Third, the numerical construction involving \(x_p=21\) and
\(K=2\pi e^{-1/42}\) is conjectural. It is a proposal for
future analysis, not a consequence of the geometric proof
or the field measurements.

These distinctions are essential to the interpretation of
the Al-Hindara framework. Conceptual unity can motivate
mathematical models, but a quantitative relation requires
operational definitions, reproducible measurements, and
a derivation or empirical test appropriate to the claim.

% ================================================================
% 7. CONCLUSION
% ================================================================
\section{Conclusion}

This paper derives the exact identity
\[
    \tan\left(\frac{\pi}{8}\right)=\sqrt{2}-1,
\]
which is the reciprocal of the conventional Silver Ratio.
Under the specified angle convention, the result gives an
exact shadow-to-height ratio for an idealized right triangle.

The five reported mango-tree observations have a mean crown
ratio of approximately \(0.41297\). The reported TOST analysis
indicates equivalence to the theoretical reference within
the stated margin of \(\pm0.03\), subject to the assumptions
and limitations of the analysis. Because the sample is
small and drawn from one location, the result is exploratory
and should not be interpreted as evidence of a universal
biological law or a causal relationship.

The paper also defines a conjectural numerical construction
using the integer \(x_p=21\). Its proposed relationship to the
Riemann zeta function remains unproved and requires explicit
mathematical definitions and independent derivation.

Finally, the relation \(Z+C+A=1\) remains a conceptual
normalization proposal until its variables and normalization
rules are formally specified. Future work should prioritize
reproducible field methods, larger and more diverse samples,
and rigorous analysis of the proposed numerical model.

% ================================================================
% REFERENCES
% ================================================================
\begin{thebibliography}{99}

\bibitem{schuirmann1987}
Schuirmann, D. J. (1987).
A comparison of the two one-sided tests procedure and the
power approach for assessing the equivalence of average
bioavailability.
\textit{Journal of Pharmacokinetics and Biopharmaceutics},
15(6), 657--680.

\bibitem{edwards1974}
Edwards, H. M. (1974).
\textit{Riemann's Zeta Function}.
Academic Press.

\bibitem{halle1978}
Hall\'e, F., Oldeman, R. A. A., and Tomlinson, P. B. (1978).
\textit{Tropical Trees and Forests: An Architectural Analysis}.
Springer-Verlag.

\end{thebibliography}

% ================================================================
% APPENDICES
% ================================================================
\appendix

\section{Supplementary Web Applications}

The following web applications are listed as supplementary
project resources. Their availability and calculations should
be independently checked before publication.

\begin{itemize}
    \item \textbf{Visual Simulation:}
    \url{https://jabri-com.vercel.app/tree/tree2.html}

    \item \textbf{TOST Field Tool:}
    \url{https://jabri-com.vercel.app/tree/tree3.html}
\end{itemize}

\section{Reproducibility Notes}

For a reproducible statistical report, the analysis archive
should include the unrounded measurements, the calculation
of each ratio, the sample standard deviation, the standard
error, the confidence interval, and the two one-sided
test statistics and p-values.

The reported summary values in this manuscript are rounded
for presentation. Calculations should be performed using
unrounded ratios derived from the original measurements.
The measurement record should also document the actual
instrument and field procedure used.

\end{document}